\documentclass[11pt,a4paper]{article}
\usepackage{fontspec}
\setmainfont{Arial}
\setsansfont{Arial}
\setmonofont{Consolas}[Scale=0.87]
\usepackage[russian]{babel}
\usepackage[margin=21mm,top=22mm,bottom=23mm,headheight=15pt]{geometry}
\usepackage{xcolor,graphicx,tabularx,booktabs,array,enumitem,fancyhdr,titlesec,caption}
\usepackage[unicode,breaklinks,hidelinks]{hyperref}
\definecolor{plum}{HTML}{642E49}
\definecolor{graphite}{HTML}{302B31}
\definecolor{muted}{HTML}{72666D}
\definecolor{warm}{HTML}{F5F0ED}
\hypersetup{pdftitle={ПУЛЬС — описание процесса создания проекта с использованием искусственного интеллекта},pdfauthor={Проект ПУЛЬС}}
\graphicspath{{assets/}}
\color{graphite}
\setlength{\parindent}{0pt}
\setlength{\parskip}{6pt}
\linespread{1.1}
\setlength{\emergencystretch}{3em}
\setlist[itemize]{leftmargin=5mm,itemsep=2pt,topsep=3pt}
\titleformat{\section}{\Large\bfseries\color{plum}}{\thesection}{0.7em}{}
\titleformat{\subsection}{\large\bfseries}{\thesubsection}{0.7em}{}
\titlespacing*{\section}{0pt}{12pt}{9pt}
\titlespacing*{\subsection}{0pt}{10pt}{5pt}
\captionsetup{font=small,labelfont=bf,skip=6pt}
\pagestyle{fancy}\fancyhf{}
\fancyhead[L]{\small\color{plum}\textbf{ПУЛЬС}}
\fancyhead[R]{\small\color{muted}Процесс создания проекта}
\fancyfoot[L]{\scriptsize\color{muted}Конкурсный прототип · вымышленные данные}
\fancyfoot[R]{\small\thepage}
\renewcommand{\headrulewidth}{0.3pt}
\newcommand{\source}[1]{{\small\color{muted}#1\par}}
\newcommand{\shot}[3]{\begin{center}\includegraphics[width=#2\linewidth,height=0.62\textheight,keepaspectratio]{#1}\captionof{figure}{#3}\end{center}}
\begin{document}
\begin{titlepage}
\vspace*{15mm}
{\fontsize{48}{52}\selectfont\bfseries\color{plum}ПУЛЬС\par}
\vspace{9mm}
{\LARGE\bfseries Описание процесса создания проекта с использованием искусственного интеллекта\par}
\vspace{8mm}
{\large Конкурс «Цифровой деканат будущего»\par}
\vspace{13mm}
\rule{\linewidth}{0.5pt}
\vspace{4mm}
Личный кабинет студента: от постановки задачи и проектных документов к работающим учебным и административным сценариям.

\vspace{5mm}
\textbf{Исходный код}\par
\url{https://github.com/GomanNik/puls-student-cabinet}

\textbf{Публичный адрес прототипа}\par
\fbox{\parbox{0.94\linewidth}{Будет добавлен после отдельной публикации.\newline На дату подготовки приложение проверено локально.}}

\textbf{Демонстрационные аккаунты}\par
\texttt{demo-a} — Дёмина Александра Ильинична, группа ПИ-24-1.\newline
\texttt{demo-b} — Лазарев Михаил Артёмович, группа ИК-23-2.\newline
Все данные вымышлены. Пароли передаются отдельно в специальном комментарии к конкурсной работе.
\vfill
{\color{muted}Материалы подготовлены 19 ноября 2026 года.\par
ПУЛЬС не является официальным сервисом университета.}
\end{titlepage}
\tableofcontents
\vspace{8mm}
\source{Документ объединяет ровно шесть тематических блоков описания процесса. Сведения о разработке взяты из ранее выполненного исследования, оригинальных заданий, проектных документов и Git. Новые изображения сняты с локального приложения. Отсутствующие исторические материалы не реконструировались под видом оригиналов.}
\clearpage
\section{Проблема студентов и выбор разделов}
Учебная информация распределена между расписанием, ведомостями, материалами, объявлениями и административной перепиской. Даже когда каждый источник доступен, студенту приходится самостоятельно связывать занятия, задания, сроки и результаты. На телефоне постоянные переходы между такими источниками особенно неудобны.

ПУЛЬС отвечает на эту проблему единым кабинетом. Это продуктовая гипотеза, сформированная из конкурсного задания и проектных требований. Опросы однокурсников и измерения экономии времени не проводились и не используются как доказательство.

\subsection{От информации к действию}
Основные разделы выбраны вокруг повседневных задач студента:
\begin{itemize}
\item \textbf{Вход и профиль} определяют личный учебный контекст.
\item \textbf{Расписание} показывает время занятий, аудитории, преподавателей и изменения.
\item \textbf{Учёба} объединяет дисциплины, посещаемость, работы, сроки, результаты и официальную аттестацию.
\item \textbf{Заявления и справки} позволяют отправить форму и следить за состоянием своей заявки.
\item \textbf{Уведомления} связывают изменение с соответствующим объектом кабинета.
\item \textbf{Обращения, материалы и календарь} дополняют эти сценарии перепиской, учебными ресурсами и общей картиной сроков. Календарь поддерживает экспорт ICS, PDF и XLSX.
\item \textbf{Поиск и ИИ-помощник} дают общую точку входа в сведения, разрешённые текущему студенту.
\end{itemize}
Разделы используют общую предметную модель: занятие связано с дисциплиной, работа — со сроком и результатом, обращение — с историей сообщений, уведомление — с событием. Данные не создаются заново для каждого экрана.

Учебная модель не сводится к сумме баллов. В разных дисциплинах могут сочетаться лекции, практические занятия, индивидуальные задания, курсовые проекты и экзамены. Отсутствие оценки не означает ноль, а незаполненное посещение не означает пропуск. Официальная подгруппа определяется отдельно для дисциплины.

\subsection{Почему университетские источники демонстрационные}
\textbf{Тестовые внешние источники — сознательная граница проекта.} Конкурсный студенческий модуль не требует разработки всей информационной системы университета и подключения к реальным персональным данным. Поэтому учебные сведения поступают через серверные адаптеры с детерминированными вымышленными наборами.

При этом собственные действия студента выполняются по-настоящему: сервер проверяет запрос, сохраняет заявление, обращение, сообщение или отметку прочтения в PostgreSQL и возвращает результат через Student API. Фоновый обработчик взаимодействует с тестовым административным источником и обновляет статусы. Сохранение операции и ответ источника остаются разными событиями.

\subsection{Архитектурная самостоятельность}
ПУЛЬС реализован как \textbf{модульный монолит}: Next.js объединяет интерфейс и собственный REST/JSON API, предметные сервисы отвечают за учебные и административные правила, Prisma — за доступ к PostgreSQL. Адаптеры университетских источников и шлюз GigaChat отделены от интерфейса. Фоновый worker выполняет доставку операций.

Такое распределение ответственности позволяет в дальнейшем заменить тестовые источники реальными интеграциями и включить студенческий модуль в более широкую сервисную экосистему. Это возможность развития архитектуры, а не утверждение, что уже создана распределённая микросервисная система. Кабинеты преподавателя, деканата и администратора в текущий объём не входят.
\source{Основания: концепция 02, требования 03, сценарии 04, архитектура 05, модель данных 06; действующие src/modules, src/server/integrations, prisma и scripts/administration-worker.ts.}
\clearpage
\section{Использованные инструменты и их роль}
\subsection{Искусственный интеллект и работа автора}
\textbf{ChatGPT} использовался для обсуждения концепции, подготовки заданий, независимого аудита документов и визуальных итераций. В истории сохранились исходные запросы на проверку и точечные исправления. Позднее через архив исходников обсуждалось объединение поиска и помощника. Полные ответы некоторых бесед недоступны, поэтому авторство каждого промежуточного изменения установить невозможно.

\textbf{Codex} выполнял работу в репозитории: создал первоначальную документационную основу, реализовывал программные блоки, исследовал причины ошибок, вносил адресные исправления и запускал проверки. В этой подготовке Codex также изменил выбор провайдера ИИ, снял реальные кадры Playwright и собрал материалы. Точная модель отдельных исторических сеансов не установлена; самостоятельная сторонняя IDE не заявляется без подтверждения.

\textbf{GigaChat} является внешним ИИ-провайдером внутри приложения. Он не указан как автор кода. 19 ноября выполнена реальная проверка OAuth и генерации моделью GigaChat-2; ответы не заменялись тестовыми. Подробности новой проверки приведены в блоке безопасности.

\subsection{Разработка, дизайн и проверка}
\begin{tabularx}{\linewidth}{@{}p{0.30\linewidth}X@{}}
\toprule\textbf{Средства}&\textbf{Фактическая роль}\\\midrule
HTML, CSS, JavaScript & Самостоятельные интерактивные макеты, сохранившиеся в архиве ревизии 09.\\
Next.js, React, TypeScript & Интерфейс, серверные обработчики и типизированная бизнес-логика.\\
Node.js, pnpm & Выполнение приложения, зависимости и команды разработки.\\
PostgreSQL, Prisma & Сессии, собственные пользовательские операции, схема и миграции БД.\\
Tailwind CSS, Radix UI, Lucide & Стилизация, интерфейсные примитивы и иконки. shadcn/ui указан в принятой UI-основе.\\
Zod, OpenAPI 3.1 & Проверка входных данных и формальные контракты API.\\
Git, GitHub, GitHub Actions & Версии исходников, публичный репозиторий и автоматизированные проверки. Автоматическое развёртывание не настроено.\\
Vitest, Playwright & Модульные, интеграционные и браузерные проверки. Playwright также создаёт актуальные скриншоты и видеозапись.\\
Chrome, CDP & Историческая диагностика интерактивности входа; Chrome используется и для текущей браузерной проверки.\\
TypeScript compiler, ESLint & Проверка типов и статический анализ кода.\\
Docker Compose, Caddy & Подготовленная схема контейнерного размещения и HTTPS. Публичное исполнение этой схемы не заявляется.\\
Argon2id, Noto Sans & Хеширование паролей и шрифт документов/выгрузок; это технические средства, а не ИИ-авторы.\\
LaTeX, Tectonic, Poppler & Исходник, сборка этого PDF и визуальная проверка его страниц.\\
\bottomrule
\end{tabularx}

Figma обсуждалась как способ дальнейшей работы с макетами, однако завершённый импорт не подтверждён. Использование Cursor, Copilot, отдельного графического редактора или конкретного генератора иллюстраций не установлено. Иллюстрации в интерфейсе присутствуют в существующих ресурсах проекта; авторство генератора по доступным источникам не определяется. Новые изображения интерфейса для этого отчёта нейросетью не создавались.
\source{Основания: исследование истории, журнал разработки 11, package.json, workflow quality.yml, архив HTML-ревизии 09 и фактические операции текущей подготовки.}
\clearpage
\section{Ключевые промпты и диалоги с ИИ}
Ниже приведены выборочные \textbf{дословные выдержки} из семи оригинальных заданий файла PULS\_SELECTED\_AI\_PROMPTS.md. Каждый пример сокращён; пропущенный текст не заменён реконструкцией. Описание результата после цитаты — редакционное изложение, а не реплика ИИ.

\subsection{П1. Документы до программирования}

\source{10.11.2026 · Codex · Выборочная цитата; остальной текст опущен.}

\begin{quote}\small На текущем этапе твоя задача — создать профессиональную систему управления проектом и документационную основу, которая впоследствии станет единственным источником требований и технических решений для всех ИИ-ассистентов, участвующих в разработке.\end{quote}

\textbf{Задача.} Задать границы прототипа, требования и общую основу для нескольких сеансов разработки.


\textbf{Результат.} Создано 27 начальных файлов. Это была основа: в ней имелись 26 требований, 14 сценариев и незаполненные разделы маршрутов OpenAPI. Пользователь затем потребовал последовательной проработки небольших групп документов.


\textbf{Проверка и предел вывода.} Факт создания и исходное содержимое подтверждены историей Codex. Количество файлов не используется как показатель готовности спецификаций.


\subsection{П2. Независимый аудит}

\source{11.11.2026, 15:10 · ChatGPT · Выборочная цитата; остальной текст опущен.}

\begin{quote}\small Не изменяй документацию. Предоставь полный доказательный отчёт и матрицы, предусмотренные инструкцией. Не называй предполагаемые ошибки подтверждёнными без прямых доказательств.\end{quote}

\textbf{Задача.} Проверить комплект по исходному архиву, отделив подтверждённые противоречия от предположений.


\textbf{Результат.} Следующее сообщение пользователя называет AUD-001–AUD-011 и отдельно H-01–H-07. Эти обозначения фиксируют объём следующего поручения.


\textbf{Проверка и предел вывода.} Оригинал запроса доступен. Полный итоговый отчёт аудита не получен; количество замечаний не выдаётся за независимо пересчитанный результат.


\clearpage

\subsection{П3. Исправить только доказанное}

\source{11.11.2026, 15:59 · ChatGPT · Выборочная цитата; остальной текст опущен.}

\begin{quote}\small Не выдавай гипотезы H-01–H-07 за ошибки. Выполни весь объём за один проход, без повторного согласования уже принятых решений.\end{quote}

\textbf{Задача.} Перепроверить доказательства, исправить только разрешённые файлы и сохранить остальные документы.


\textbf{Результат.} В задании требовались ZIP изменённых файлов, машинные проверки и доказательство неизменности остального комплекта.


\textbf{Проверка и предел вывода.} Подтверждены постановка и критерии. Итоговый ZIP и его контрольные суммы недоступны; закрытие всех 11 пунктов не утверждается.


\subsection{П4. Безопасность при выходе из кабинета}

\source{13.11.2026 · Codex · Выборочная цитата; остальной текст опущен.}

\begin{quote}\small Не допускать commit новой P-мутации после завершившегося отзыва той же сессии.\end{quote}

\textbf{Задача.} Устранить гонку между выходом и записью заявления, обращения, сообщения или прочтения уведомления.


\textbf{Результат.} В матрице исправлений зафиксирована повторная проверка сессии с блокировкой внутри той же транзакции. Сетевые вызовы отделены от транзакции.


\textbf{Проверка и предел вывода.} Исторический этап разрешал реализацию без запуска тестов. Статическая трассировка 77 требований, 31 сценария, 20 наборов и 49 HTTP-операций не называется испытанием приложения.


\clearpage

\subsection{П5. Найти причину неработающего входа}

\source{14.11.2026, 09:07 · Codex · Выборочная цитата; остальной текст опущен.}

\begin{quote}\small Не ограничивайся чтением исходников: установи, на каком этапе перестаёт работать клиентская часть.\end{quote}

\textbf{Задача.} Исправить форму, которая отображалась, но не принимала ввод. Предыдущие изменения формы и CSP не решили проблему.


\textbf{Результат.} Диагностика связала незавершённую гидрацию с тем, что локальный HTTPS-прокси не передавал WebSocket Upgrade. Добавлена ограниченная обработка соединения Next.js.


\textbf{Проверка и предел вывода.} В журнале зафиксированы ответ 101, интерактивность полей, ошибочный вход и вход обоих студентов. Исторического изображения самого дефекта нет; его инсценировка не создавалась.


\subsection{П6. Пересмотреть неудачную компоновку}

\source{18.11.2026, 13:52 · Codex · Выборочная цитата; остальной текст опущен.}

\begin{quote}\small **Важное изменение требований:** предыдущее указание использовать полноширинные вертикальные списки отменяется. Оно оказалось неудачным на широких экранах.\end{quote}

\textbf{Задача.} Сделать записи работ компактными, убрать дублирование результатов и сохранить серверные фильтры и переходы.


\textbf{Результат.} Пользователь задал сетку до трёх карточек, две на среднем экране и одну на мобильном. Неполный ряд не должен растягиваться; отсутствие результата нельзя заполнять вымышленной оценкой.


\textbf{Проверка и предел вывода.} Исходная команда подтверждена. Текущая страница снята с работающего приложения; она не заменяет отсутствующий исторический кадр неудачной версии.


\clearpage

\subsection{П7. Сохранять прочтение уведомлений}

\source{18.11.2026, 22:18 · Codex · Выборочная цитата; остальной текст опущен.}

\begin{quote}\small 5. Если обновление на сервере завершилось ошибкой, не показывай ложное подтверждение прочтения.\end{quote}

\textbf{Задача.} Связать фактическое открытие уведомления с API, обновлением счётчика и сохранением прочтения после перезагрузки.


\textbf{Результат.} Исправление выделено в коммит de1fb52; добавлен адресный браузерный регрессионный тест. Позднее пользователь уточнил, что счётчик остаётся только в верхней панели.


\textbf{Проверка и предел вывода.} Подтверждены Git-коммит и успешный CI этого SHA. CI не включал весь браузерный набор и не доказывает публичную доступность.


\source{Исторический источник: PULS\_SELECTED\_AI\_PROMPTS.md, П1–П7. Git-проверка: \url{https://github.com/GomanNik/puls-student-cabinet/actions/runs/44352118677}. Публичные материалы не содержат приватных ссылок на беседы.}
\clearpage
\section{Скриншоты ключевых этапов}
\subsection{Сохранившиеся визуальные варианты}
Идея единого кабинета подтверждена первоначальным заданием и концепцией. Отдельного изображения самой первой идеи среди проверенных материалов нет. Ниже показаны доступные архивные HTML-превью, а не якобы восстановленная первая версия.

\shot{revision09_index_1440.png}{1}{Главная страница HTML-макета ревизии 09, 1440 px. Архивный ориентир композиции; не доказательство работы API или БД.}
\source{Источник: PULS\_HTML\_revision\_09.zip, previews/revision09\_index\_1440.png. Изображение перенесено без редактирования.}
\clearpage
\shot{notifications_1440.png}{0.86}{Архивное превью уведомлений, сохранённое в комплекте ревизии 09.}
\shot{revision09_notifications_1440.png}{0.86}{Вариант уведомлений с обозначением revision09 в том же архиве.}
\source{Оба изображения подлинные. Их имена и визуальное различие подтверждены; точный порядок всех промежуточных версий не установлен. Неудачное решение полноширинных списков документировано в П6, но отдельного исторического кадра этого дефекта нет.}
\clearpage
\subsection{Работающий кабинет: вход и персональный контекст}
Следующие изображения сняты 19.11.2026 через Playwright с локального \texttt{https://localhost:3443}. Настольный viewport — 1440 × 960 px; мобильный — 360 × 800 px. Это реальные кадры браузера, без подмены содержимого страницы. Локальная съёмка не означает публикацию приложения.
\shot{login-desktop.png}{0.90}{Форма входа. Пароли и токены на кадре отсутствуют.}
\shot{demo-a-dashboard-1440.png}{0.90}{Главная страница первого студента после входа.}
\clearpage
\shot{demo-a-profile-1440.png}{0.90}{Профиль demo-a: Александра Дёмина, ПИ-24-1.}
\shot{demo-b-profile-1440.png}{0.90}{Профиль demo-b: Михаил Лазарев, ИК-23-2. Данные обоих студентов вымышлены.}
\clearpage
\subsection{Учебные и административные сценарии}
\shot{demo-a-schedule-1440.png}{0.90}{Расписание после выбора своей группы.}
\shot{discipline-results-1440.png}{0.90}{Дисциплина и результаты учебной деятельности.}
\clearpage
\shot{request-created-1440.png}{0.90}{Заявка, отправленная через форму: собственная запись и история состояния.}
\shot{demo-a-notifications-1440.png}{0.90}{Единые уведомления с переходами к объектам.}
\clearpage
\shot{appeal-thread-1440.png}{0.90}{Обращение с историей переписки.}
\shot{gigachat-1440.png}{0.90}{Ответ настоящего GigaChat через Student API. Маркер провайдера отображён в интерфейсе.}
\clearpage
\subsection{Мобильный интерфейс}
\begin{center}
\begin{minipage}[t]{0.46\linewidth}\centering
\includegraphics[width=0.83\linewidth]{mobile-dashboard-360.png}
\captionof{figure}{Главная, 360 px.}
\end{minipage}\hfill
\begin{minipage}[t]{0.46\linewidth}\centering
\includegraphics[width=0.83\linewidth]{mobile-schedule-360.png}
\captionof{figure}{Расписание, 360 px.}
\end{minipage}
\vspace{7mm}
\begin{minipage}[t]{0.46\linewidth}\centering
\includegraphics[width=0.83\linewidth]{mobile-request-360.png}
\captionof{figure}{Заявка, 360 px.}
\end{minipage}\hfill
\begin{minipage}[t]{0.46\linewidth}\centering
\includegraphics[width=0.83\linewidth]{mobile-notifications-360.png}
\captionof{figure}{Уведомления, 360 px.}
\end{minipage}
\end{center}
\source{Архивный мобильный макет имел ширину 390 px и не используется как доказательство требования 360 px. Кадры выше получены отдельно с актуального приложения. Исторических скриншотов всех диалогов ChatGPT и Codex нет; подтверждённые оригинальные тексты приведены в предыдущем блоке.}
\clearpage
\section{Вход, хранение данных и безопасность}
\subsection{От пароля к серверной сессии}
Студент вводит логин и пароль. Сервер проверяет пароль по хешу Argon2id; исходный пароль в БД не сохраняется. Число неудачных попыток ограничивается. После входа браузер получает непрозрачный идентификатор сессии, а сервер хранит её состояние в PostgreSQL.

Cookie сессии имеет флаги Secure, HttpOnly и SameSite=Lax: она передаётся по HTTPS и недоступна обычному клиентскому JavaScript. Выход отзывает сессию на сервере. Для операций записи дополнительно проверяются защищённый канал, Origin и действующая сессия.

\subsection{Почему нельзя открыть чужую запись}
Права проверяются на сервере, а не наличием кнопки в интерфейсе. Собственные заявки, обращения и сообщения выбираются с учётом владельца. Учебная область строится из подтверждённой идентичности студента, периода и официальных назначений. Подстановка чужого идентификатора в URL не меняет владельца текущей сессии.

Два демонстрационных аккаунта имеют разные профили и учебный контекст. При текущей подготовке вход проверен в отдельных браузерных контекстах. Адресные проверки подстановки профиля, чужой заявки и обращения описываются по фактическим результатам, без вывода о гарантированном отсутствии любых будущих уязвимостей.

\subsection{Собственные операции и внешние сведения}
Заявления, обращения, сообщения и индивидуальное прочтение уведомлений сохраняются в реальной PostgreSQL. Сведения об учебной деятельности приходят из серверных тестовых источников и не превращаются в независимо редактируемые копии для каждого раздела. Студент не может изменить официальную оценку, посещаемость, расписание или подгруппу.

Доставка административной операции и её предметный результат различаются. Сохранённая заявка не объявляется одобренной до получения соответствующего ответа источника. Worker обрабатывает очередь; повторные отправки и изменения защищены предусмотренными контрактами и транзакциями.

\subsection{Настоящий ИИ и серверные секреты}
В ходе финальной подготовки обнаружено, что клиентская константа \texttt{DEMO\_MODE = true} возвращала локальный ответ, обходя серверный адаптер. Она удалена. Интерфейс обращается к Student API; основной серверный провайдер по умолчанию — \texttt{GIGACHAT}. Тестовый \texttt{DEMO\_MOCK} включается явно в автоматизированных проверках. Ошибка настоящего провайдера не подменяется демонстрационным ответом.

Авторизационный ключ остаётся в защищённой серверной конфигурации. Браузер не получает ключ или OAuth-токен. Серверный адаптер использует OAuth, ограничение времени ожидания, проверку структуры ответа, разрешённые факты текущего студента и ограниченные повторные запросы. Неверный ключ, ограничение запросов, неизвестная модель и недоступность сервиса различаются.

Для GigaChat проверка TLS сохранена; Node.js запущен с дополнительным доверенным набором сертификатов. В реальном тесте 19.11.2026 OAuth вернул HTTP 200, три запроса генерации — HTTP 200, неверная модель — 404, заведомо неверный ключ — 401. Проверено продолжение диалога. Реквизиты, токены и сырые заголовки в конкурсные материалы не включены.

Текущие адресные проверки: 13 модульных тестов адаптера и выбора провайдера, 6 интеграционных тестов разрешённого контекста, проверка типов и ESLint. Реальный сетевой тест выполнялся отдельно от mock-проверок. Публичный HTTPS и доступность из России проверяются после отдельного развёртывания; локальная проверка их не заменяет.
\source{Реализация: модули auth, student-context, operations, community и assistant; server/integrations/assistant-provider.ts. Протокол текущей проверки хранится локально без секретов. Описание OAuth сверено с \url{https://developers.sber.ru/docs/ru/gigachat/api/reference/rest/gigachat-api}.}
\clearpage
\section{Этапы, выполненные вручную}
\subsection{Постановка и принятие решений}
Автор определял границы студенческого модуля, отделял обязательные функции от расширений и задавал правила демонстрационных данных. Он принимал предметные и визуальные решения, проверял предложения ИИ по конкурсному заданию и уточнял задания исполнителям. ChatGPT и Codex не являлись владельцами продукта.

10 ноября автор отказался от массовой подготовки всего комплекта документов и потребовал последовательной проработки. В заданиях на аудит он закрепил требование доказательств, отделение гипотез от ошибок и запрет переписывать неподтверждённые части проекта. Организация работы изменилась по результатам оценки первых материалов.

В дизайне автор выбирал сдержанную всесезонную палитру и функциональную компоновку, оценивал варианты по реальным изображениям и возвращал неудачные решения на доработку. Отказ от полноширинных списков в пользу карточек — подтверждённый пример пересмотра ранее принятого указания после проверки на широком экране.

\subsection{Контроль качества и корректировка ИИ}
Автор обнаруживал несоответствия при использовании приложения: неработающую интерактивность входа, некорректное количество занятий, проблемы статусов и прочтения уведомлений. Он требовал искать первопричину, а не скрывать симптом, сохранять архитектуру и не ослаблять защиту ради демонстрации.

Работа распределялась по результату: ChatGPT помогал обсуждать требования, формулировать проверяемые задания и оценивать варианты; Codex работал с файлами, кодом, командами и браузером; автор выбирал и принимал решения, проверял смысл и корректировал направление. Такое разделение не означает, что каждый документ написал ChatGPT, а весь код — только Codex: авторство конкретной итерации определяется источником.

При подготовке этой версии автор отдельно потребовал основной GigaChat, реальные снимки работающего кабинета, PDF из LaTeX и отсутствие фиктивной ссылки на публикацию. Он предоставил доступы к двум тестовым аккаунтам. Исполнитель сохранил их локально вне Git и использовал для браузерной проверки.

\subsection{Граница подтверждённого результата}
История восстановлена по доступным первоисточникам и существующим артефактам. Полный экспорт всех бесед, все промежуточные дизайн-версии и результаты первоначального аудита недоступны. Поэтому неизвестные действия не приписываются автору или ИИ, а отсутствующие изображения не имитируются.

На дату документа приложение демонстрируется локально; публикация отложена отдельным решением автора. До отправки необходимо добавить настоящий публичный адрес, проверить доступность из России и обеспечить работу до объявления итогов 28.11.2026. Пароли передаются в специальном комментарии, отдельно от PDF, исходников и видео.
\source{Основания: ранее выполненное исследование PULS\_AI\_DEVELOPMENT\_HISTORY.md; семь оригинальных заданий; журнал разработки; фактическая текущая подготовка. Частные ссылки и локальные секреты из этих источников в отчёт не перенесены.}
\end{document}
